\PuzzleChapterIntro{The Three Bead Readings}{Geometry \textbullet\ Vault Mechanics \textbullet\ Ratios}{This puzzle uses vector geometry and locus definitions. You must convert "secant ratios" into distance equations to find the locus of the point $P$, and then use the "midpoint distances" to calculate the size of that locus.}
\Lore{
In the Astral Annex, diviners read geometry the way accountants read ledgers: by invariants.

A sphere of radius \(R=9\) is suspended on three chains. Four beads \(W,X,Y,Z\) are fixed to the glass, and a fifth point \(P\) may wander.
The diviners do not measure angles directly. They measure \emph{lengths} and \emph{powers}, then declare whether a reading can be made.

Three bead-readings are written in the margin: the squared distances between the midpoints of opposite bead-pairs.
The rest, they claim, can be deduced.
}

\Problem

The vault’s \emph{heart} is the point \(H\) that is equally distant from every point of the wall; that common distance is \(9\).

A lantern is held at a point \(P\) strictly inside the vault. For each nail (say \(W\)), extend the straight line through \(W\) and \(P\) until it meets the wall again on the far side; call that second intersection \(W_1\neq W\). Define the \emph{weighing} of nail \(W\) to be the ratio
\[
\frac{WP}{PW_1},
\]
where both lengths are measured along the same straight line \(W\!-\!P\!-\!W_1\).
Thus the four weighings are \(\frac{WP}{PW_1},\frac{XP}{PX_1},\frac{YP}{PY_1},\frac{ZP}{PZ_1}\).

The brass rule says the sum of the four weighings equals \(4\):
\[
\frac{WP}{PW_1}+\frac{XP}{PX_1}+\frac{YP}{PY_1}+\frac{ZP}{PZ_1}=4.
\]

Separately, the ledger describes the keeper’s \emph{pairing ritual}:

\begin{itemize}
\item On a night, the keeper connects two disjoint pairs of nails with taut wire (so every nail is used exactly once that night).
\item The keeper presses a wax bead to the midpoint of each wire.
\item The keeper measures the straight distance between the two beads.
\item The keeper repeats until, across all nights, \emph{every pair of nails has been connected exactly once}.
\end{itemize}

The three recorded bead-distances (in no particular order) are
\[
12,\quad 8,\quad 4.
\]

The Archive claims that the set of all lantern-positions $P$ obeying the brass rule forms one perfectly round \emph{inner shell}.
\VaultDirective{What is the \textbf{diameter length} of that inner shell?}

\SolutionPage
\small
Let \(H\) be the sphere’s center (\(R=9\)). For any bead \(A\in\{W,X,Y,Z\}\), the secant through \(P\) meets the sphere again at \(A_1\), and
\[
P A\cdot P A_1=R^2-HP^2 \quad\Rightarrow\quad \frac{PA}{PA_1}=\frac{PA^2}{R^2-HP^2}.
\]
So the Oracle equation
\(\sum_{A}\frac{PA}{PA_1}=4\) is equivalent to
\[
\sum_A PA^2=4(R^2-HP^2). \tag{1}
\]

Let \(C\) be the centroid of \(WXYZ\). The identity \(\sum_A \overrightarrow{PA} = 4\overrightarrow{PC}\) implies the quadratic decomposition
\[
\sum_A PA^2 = 4\,PC^2+\sum_A CA^2. \tag{2}
\]
Apply (2) at \(P=H\): since \(HA=R\) for all beads,
\[
4R^2=\sum_A HA^2 = 4\,HC^2+\sum_A CA^2 \ \Rightarrow\ \sum_A CA^2=4(R^2-HC^2). \tag{3}
\]
Substitute (2) and (3) into (1):
\[
4PC^2+4(R^2-HC^2)=4(R^2-HP^2)
\ \Longrightarrow\ 
HP^2+PC^2=HC^2.
\]
Thus the locus of \(P\) is exactly the sphere whose diameter is \(HC\). So the asked diameter is \(HC\).

Let \(a_W,a_X,a_Y,a_Z\) be the position vectors of the beads and \(c=\frac{a_W+a_X+a_Y+a_Z}{4}\) (\(|c|=HC\)).
Then
\[
\sum_{i<j}\|a_i-a_j\|^2=4\sum_i\|a_i\|^2-\Bigl\|\sum_i a_i\Bigr\|^2
=16R^2-16HC^2=16(R^2-HC^2). \tag{4}
\]
For midpoints \(M_{WX},M_{YZ}\) etc.,
\[
\|M_{WX}-M_{YZ}\|^2=\tfrac14\|(W+X)-(Y+Z)\|^2,
\]
and summing the three “opposite-midpoint” squares gives
\[
d_1^2+d_2^2+d_3^2=\tfrac14\sum_{i<j}\|a_i-a_j\|^2. \tag{5}
\]
Given $d_1=12,d_2=8,d_3=4$, we have $d_1^2+d_2^2+d_3^2=12^2+8^2+4^2=224$, so by (5) and (4):
\[
224=\tfrac14\cdot 16(R^2-HC^2)\ \Rightarrow\ R^2-HC^2=56.
\]
With \(R^2=81\), we get \(HC^2=25\), so \(HC=5\). Therefore the locus sphere has radius \(HC/2=\boxed{\tfrac52}\), so its \emph{diameter} is \(\boxed{5}\).

\FinalAnswer{5}

\NextPage
